\section{示例}

\begin{frame}[fragile]{文档类}

	首先，使用\texttt{beamer}文档类配合\texttt{ust}主题。需要注意的是，\texttt{ust.sty}文件应该与\texttt{main.tex}文件放在同一目录中。
	
\begin{block}{关于文档类的导言}
\begin{lstlisting}[language=TeX]
\documentclass[10pt,aspectratio=169]{beamer}
\usepackage{ust}
\end{lstlisting}
\end{block}
	
	您可以将\texttt{aspectratio}改为\texttt{43}来调整幻灯片的宽高比为4:3。
\end{frame}

\begin{frame}[fragile]{元数据}
	您可以更改标题页上显示的元数据：

\begin{block}{元数据}
\begin{lstlisting}[language=TeX]
\title{您的标题}
\subtitle{您的副标题}
\author{第一作者，第二作者}
\institute[ECE]{电子与计算机工程学系}
\date{日期}
\end{lstlisting}
\end{block}

	设置完成后，您可以在文档正文中使用\verb|\maketitle|命令来生成标题页。
\end{frame}

\begin{frame}[fragile]{文本}{按序列显示}

	\begin{itemize}[<+->]
		\item 典型的幻灯片都有项目符号。
		\item 这些内容可以按顺序显示。
		\item 渲染时，它们会被分成多张幻灯片。
	\end{itemize}

\end{frame}

\begin{frame}{图像}

	添加图像的方法与在普通\LaTeX 中相同：
	
	\begin{figure}[hbt]
		\includegraphics[width=0.5\textwidth]{assets/template/logo-name.pdf}
 		\caption{图像示例}
	\end{figure}
	
\end{frame}

\begin{frame}{分栏}
	分割页面是一种简单而常见的做法。

	通常，一边放图片，另一边放文字：
	
	\vspace{20pt}
	
	\begin{columns}
		\begin{column}{0.6\textwidth}
			这是第一栏。\\[10pt]
			您可以在这里放一些文字。
		\end{column}
		
		\begin{column}{0.3\textwidth}
			这是第二栏。\\[10pt]
			您可以在这里放一些图片或表格。
		\end{column}
	\end{columns}

\end{frame}

\begin{frame}{字体}
	
	选择字体时的优先原则是可读性。
	
	这里我们给出一些建议：
	\begin{itemize}
		\item 只有在高分辨率投影仪或显示器上才使用衬线字体（默认）。
		\item 其他情况下使用\textsf{无衬线字体}。
		\item 使用\textit{斜体}或\textbf{粗体}字体来强调或突出重点。
		\item 我们还提供了用于强调的\alert{警示}字体。
		\item 使用\texttt{等宽字体}来显示代码。
	\end{itemize}

\end{frame}

\begin{frame}{参考文献}

	您可以在\texttt{.bib}文件中添加参考文献，并使用\texttt{biblatex}包来引用它们\cite{liuAndyshiyiliuHKUSTBeamerTemplate2025}。

\end{frame}